\section{Recipe 4: Tạo ảnh nghệ thuật từ văn bản với Stable Diffusion}
\label{sec:recipe_stable_diffusion}

Stable Diffusion là mô hình sinh ảnh từ văn bản (text-to-image) mạnh mẽ và mã nguồn mở. Không giống các API đóng, bạn có thể chạy nó hoàn toàn cục bộ trên GPU của mình. Recipe này hướng dẫn ba kịch bản sử dụng: sinh ảnh từ văn bản, biến đổi ảnh có sẵn, và kiểm soát chính xác bố cục với ControlNet.

\subsection{Bước 1: Sinh ảnh từ văn bản (Text-to-Image)}
\label{subsec:sd_text2img}

Thư viện \texttt{diffusers} của Hugging Face là cách dễ nhất để sử dụng Stable Diffusion. Pipeline encapsulate toàn bộ quá trình: mã hóa text prompt, khử nhiễu latent, và giải mã thành ảnh.

\begin{minted}[breaklines=true, fontsize=\small]{python}
from diffusers import StableDiffusionPipeline, DPMSolverMultistepScheduler
import torch

model_id = "stabilityai/stable-diffusion-2-1"

# Tai pipeline voi float16 de giam VRAM (can ~5GB VRAM)
pipe = StableDiffusionPipeline.from_pretrained(
    model_id,
    torch_dtype=torch.float16
)

# Thay scheduler mac dinh bang DPM-Solver++: nhanh hon, chat luong tuong duong
pipe.scheduler = DPMSolverMultistepScheduler.from_config(pipe.scheduler.config)
pipe = pipe.to("cuda")

# Kich hoat attention slicing: giam VRAM them ~1GB, cham hon mot chut
pipe.enable_attention_slicing()

# Sinh anh
image = pipe(
    prompt="a photorealistic image of a cat wearing a wizard hat, "
           "dramatic lighting, 4k, highly detailed",
    negative_prompt="blurry, low quality, distorted, watermark, cartoon",
    num_inference_steps=25,    # Nhieu buoc = chat luong cao hon, nhung cham hon
    guidance_scale=7.5,        # CFG scale: cao = bam sat prompt, thap = sang tao hon
    height=512,
    width=512,
    generator=torch.Generator("cuda").manual_seed(42),  # Ket qua tai hien duoc
).images[0]

image.save("wizard_cat.png")
print("Da luu anh tai wizard_cat.png")
\end{minted}

\begin{mdframed}[style=infobox]
\textbf{Nghệ thuật viết prompt hiệu quả}

Prompt tốt thường có cấu trúc: \textit{[chủ thể]} + \textit{[phong cách]} + \textit{[ánh sáng/màu sắc]} + \textit{[từ khóa chất lượng]}.

\textbf{Ví dụ:} \textit{"portrait of a Vietnamese woman, traditional ao dai, golden hour lighting, oil painting, highly detailed, artstation"}.

\textbf{Negative prompt} quan trọng không kém: \textit{"blurry, low quality, distorted, extra fingers, watermark, text, bad anatomy"}.

Một số từ khóa tăng chất lượng phổ biến: \texttt{4k}, \texttt{highly detailed}, \texttt{photorealistic}, \texttt{masterpiece}, \texttt{trending on artstation}, \texttt{sharp focus}. Ngược lại, \texttt{guidance\_scale} thấp (5--6) cho ảnh sáng tạo hơn nhưng ít bám prompt; cao (8--12) bám prompt nhưng đôi khi cứng nhắc.
\end{mdframed}

\subsection{Bước 2: Image-to-Image (Img2Img)}
\label{subsec:sd_img2img}

Img2Img cho phép bạn bắt đầu từ một ảnh có sẵn (sketch, ảnh chụp) và biến đổi nó theo phong cách mới thông qua prompt. Tham số \texttt{strength} kiểm soát mức độ biến đổi: 0 giữ nguyên hoàn toàn ảnh gốc, 1 bỏ qua hoàn toàn.

\begin{minted}[breaklines=true, fontsize=\small]{python}
from diffusers import StableDiffusionImg2ImgPipeline
from PIL import Image

img2img_pipe = StableDiffusionImg2ImgPipeline.from_pretrained(
    model_id,
    torch_dtype=torch.float16
)
img2img_pipe = img2img_pipe.to("cuda")

# Anh dau vao: co the la sketch, anh thuc, hoac anh thu cong
init_image = Image.open("sketch.jpg").resize((512, 512))

result = img2img_pipe(
    prompt="a beautiful oil painting of a mountain landscape, "
           "warm colors, impressionist style",
    image=init_image,
    strength=0.75,          # 0.0 = giu nguyen anh goc, 1.0 = sinh hoan toan moi
    guidance_scale=7.5,
    num_inference_steps=30,
).images[0]

result.save("mountain_painting.png")
\end{minted}

\subsection{Bước 3: ControlNet --- Kiểm soát bố cục chính xác}
\label{subsec:sd_controlnet}

ControlNet cho phép điều khiển cấu trúc của ảnh được sinh ra thông qua một "ảnh điều kiện" (condition image): cạnh Canny, skeleton người, bản đồ chiều sâu, v.v. Đây là công cụ mạnh khi bạn cần giữ nguyên bố cục nhưng thay đổi phong cách.

\begin{minted}[breaklines=true, fontsize=\small]{python}
import cv2
import numpy as np
from diffusers import ControlNetModel, StableDiffusionControlNetPipeline
from PIL import Image

def get_canny_edges(image_path, low_threshold=100, high_threshold=200):
    """Trich xuat canh Canny tu anh de dung lam dieu kien ControlNet."""
    image = cv2.imread(image_path)
    image_gray = cv2.cvtColor(image, cv2.COLOR_BGR2GRAY)
    edges = cv2.Canny(image_gray, low_threshold, high_threshold)
    # Chuyen ve anh RGB 3 kenh
    edges_rgb = np.stack([edges, edges, edges], axis=-1)
    return Image.fromarray(edges_rgb)

# Tai ControlNet pretrained cho Canny edges
controlnet = ControlNetModel.from_pretrained(
    "lllyasviel/sd-controlnet-canny",
    torch_dtype=torch.float16
)

pipe = StableDiffusionControlNetPipeline.from_pretrained(
    "runwayml/stable-diffusion-v1-5",
    controlnet=controlnet,
    torch_dtype=torch.float16,
)
pipe = pipe.to("cuda")
pipe.enable_attention_slicing()

# Lay ban do canh tu anh nguon
source_image = "architecture.jpg"
canny_image = get_canny_edges(source_image)

# Sinh anh moi giu nguyen cau truc cua anh nguon
result = pipe(
    prompt="a watercolor painting of a building, soft colors, artistic",
    image=canny_image,          # ControlNet dieu kien
    num_inference_steps=20,
    guidance_scale=7.5,
).images[0]

result.save("controlled_result.png")
\end{minted}

\subsection{Các mô hình Stable Diffusion phổ biến}
\label{subsec:sd_models_table}

\begin{table}[h]
\centering
\caption{Một số mô hình Stable Diffusion phổ biến và trường hợp sử dụng}
\label{tab:sd_models}
\begin{tabular}{lp{7cm}}
\hline
\textbf{Mô hình} & \textbf{Điểm mạnh và usecase} \\
\hline
SD v2.1            & Mục đích chung, ảnh thực, phong cảnh \\
SD XL (SDXL)       & Độ phân giải cao (1024px), chi tiết tốt hơn \\
SDXL-Turbo         & Sinh ảnh 1--4 bước, real-time interactive \\
Realistic Vision   & Ảnh chân thực, portrait, ảnh chụp \\
DreamShaper        & Kết hợp nghệ thuật và thực tế, linh hoạt \\
Deliberate         & Nghệ thuật chi tiết, nhân vật, concept art \\
\hline
\end{tabular}
\end{table}
